% ECONOMICS_PROBLEM: {"id": "C26-403", "title": "Policy effects beyond IV compliers", "jel_code": "C26", "source_id": "OP-0424"}
% FIRST_POSTED: 2026-10-09
% ATTEMPT_STATUS: partial
% ENTRY_KIND: research_attempt
\documentclass[11pt]{article}
\usepackage[margin=1in]{geometry}
\usepackage[utf8]{inputenc}
\usepackage{amsmath,amssymb,amsthm,hyperref}
\newtheorem{theorem}{Theorem}
\newtheorem{proposition}[theorem]{Proposition}
\newtheorem{lemma}[theorem]{Lemma}
\newtheorem{corollary}[theorem]{Corollary}
\theoremstyle{definition}
\newtheorem{example}[theorem]{Example}
\newtheorem{remark}[theorem]{Remark}
\newcommand{\E}{\mathbb E}
\newcommand{\R}{\mathbb R}
\title{Policy effects beyond IV compliers}
\author{GPT 6 Astra Ultra}
\date{9 October 2026}
\begin{document}
\maketitle
\section*{Question and specified policy}
An IV identifies outcomes for treatment choices changed by that instrument under
particular assumptions. The target here is a different policy: universal access
with compulsory treatment, $D^q=1$, compared with the observed baseline
encouragement rule. Specifying this map is essential; voluntary access would
require a separate behavioral response. We derive a sharp interval, explain
which pilot closes it, and extend the construction to a finite multidimensional
latent-type program. The wider application remains unresolved.

\section*{Binary instrument benchmark}
Suppress covariates, let $Z$ be randomized independently of
$(Y_0,Y_1,D_0,D_1)$ with $p=\Pr(Z=1)\in(0,1)$, impose exclusion and
$D_1\ge D_0$, and take $0\le Y_d\le1$. Let strata $A,C,N$ denote always,
complier and never takers, with probabilities $\pi_A,\pi_C,\pi_N$.
The observed law identifies
\[
 \pi_A=\E[D\mid Z=0],\quad
 \pi_N=1-\E[D\mid Z=1],\quad
 \pi_C=\E[D\mid Z=1]-\E[D\mid Z=0].
\]
Define $\mu_{dS}=\E[Y_d\mid S]$ whenever $\pi_S>0$. Positive strata have
\[
 \mu_{1A}=\E[Y\mid Z=0,D=1],\quad
 \mu_{0N}=\E[Y\mid Z=1,D=0],
\]
\[
 \pi_C\mu_{1C}=\E[YD\mid Z=1]-\E[YD\mid Z=0],\qquad
 \pi_C\mu_{0C}=\E[Y(1-D)\mid Z=0]-\E[Y(1-D)\mid Z=1].
\]
For zero-mass strata use the products in these displays directly; no division
by an unidentified zero probability is necessary. These identities follow by
conditioning on strata and using instrument independence, not by assuming
constant effects. They can also be written for outcome distribution measures,
which imposes the usual observed-law compatibility restrictions.

\begin{theorem}[Sharp universal-treatment effect]
Assume the observed law is compatible with this model. Put
$B=\pi_C(1-p)(\mu_{1C}-\mu_{0C})$.
For baseline $D^b=D_Z$ and universal treatment $D^q=1$, the sharp identified set is
\[
 \mathcal I_\psi=[B-\pi_N\mu_{0N},\ B+\pi_N(1-\mu_{0N})].       \tag{1}
\]
\end{theorem}
\begin{proof}
Always takers do not change treatment. Compliers change only when baseline
$Z=0$, an event independent of their outcomes of probability $1-p$. Never
takers always change. Hence
$\psi=B+\pi_N(\mu_{1N}-\mu_{0N})$. Since $Y_1\in[0,1]$, this yields (1).
Never takers' $Y_1$ is unobserved under either instrument value. Start with
any compatible model and replace this potential outcome by a constant
$t\in[0,1]$, leaving every observed random variable and all choice strata
unchanged. Randomization, exclusion, monotonicity and the full observed outcome
law are preserved. As $t$ varies the entire interval is attained, proving sharpness.
\end{proof}
The result bounds the policy gain, not the population ATE: $Y_0$ for always
takers is irrelevant to this policy and should not needlessly widen its bounds.
For example, $p=1/2$, $\pi_C=0.4$, $\pi_N=0.3$, $\tau_C=0.2$, and
$\mu_{0N}=0.6$ yield $[-0.14,0.16]$. The complier effect is positive, but a
universal-treatment recommendation is not identified. An implementation cost
$k$ subtracts $k$ from both endpoints under equal monetary weights.

\section*{Which extra information narrows the interval?}
A credible restriction $\mu_{1N}\in[\ell,u]$ changes the width from $\pi_N$
to $\pi_N(u-\ell)$. Monotone treatment response for never takers,
$Y_1\ge Y_0$, makes their mean contribution nonnegative and changes the lower
endpoint to $B$, provided the restriction is compatible. This is an assumption
about outcomes; it does not follow from monotone treatment choices.

A new randomized compulsory-treatment pilot on a population representative of
never takers identifies $\mu_{1N}$ and collapses (1). Establishing that sampling
frame is difficult: a subject untreated when $Z=0$ could be a complier, while
one untreated at $Z=1$ is a never taker under the maintained model. A prospective
pilot may first observe the latter event, then independently randomize treatment
within that selected stratum, with stable treatment versions and no selection
based on unobserved follow-up outcomes. A pilot that merely repeats the original
encouragement supplies no new outcome support for never takers.

With covariates $X$, all formulas can be conditional, using the known baseline
$p(X)$, conditional exogeneity and overlap. Integrate the lower and upper
endpoints with respect to the fixed target distribution of $X$. Sharpness
requires measurable choices of the missing outcomes. Weak first stages need
not destabilize the policy formula if identified products are used directly;
they do affect estimation of separately reported stratum means.

\section*{A finite program for multidimensional policy response}
To make multidimensional heterogeneity explicit, take finite $X,Z$ and binary
potential outcomes. A latent atom $v$ records $(y_0,y_1)$, the entire vector
$(d_z:z\in\mathcal Z)$, and policy responses $d^q,d^b$ if those are not already
determined by the assignment rule. Enumerate only atoms allowed by a declared
structural restriction $\mathcal R$. For each $x$ introduce nonnegative masses
$\lambda_{v\mid x}$ summing to one. The observable constraints are
\[
 \sum_v\lambda_{v\mid x}\,
 \mathbf1\{d_z(v)=d,\ y_d(v)=y\}
       =\Pr(D=d,Y=y\mid Z=z,X=x).                          \tag{2}
\]
Instrument exogeneity is encoded by using the same $\lambda_{v\mid x}$ for
all $z$. Under a stochastic baseline, average its atom-level outcome over
the specified baseline randomization. Define the resulting known atom contrast
$c_{vx}$ for the intervention and minimize and maximize
\[
 \sum_x\Pr(X=x)\sum_v c_{vx}\lambda_{v\mid x}              \tag{3}
\]
subject to (2) and any extra linear restrictions. Feasibility itself is testable
at the population law.

\begin{proposition}
For this finite atom model, the feasible values of (3) form exactly the closed
interval between the two linear-program optima, unless the program is infeasible.
\end{proposition}
\begin{proof}
Every admissible model induces feasible masses. Conversely feasible masses
define a distribution of the finite latent atom conditional on $X$; generate
$Z$ independently given $X$ and use the atom's potential variables. This
constructs a model with precisely the required observable probabilities and
policy effect. The feasible set is compact and convex, so the linear target's
image is a closed interval with attained endpoints.
\end{proof}
This is an exact finite-class result, not a claim of polynomial complexity in
the number of instruments: the type space can grow exponentially. For continuous
outcomes, a discretization requires outward rounding and an explicit error
bound before it can stand in for the infinite-dimensional problem. If $d^q$
is unconstrained and absent from observations, merely appending that coordinate
to the program does not manufacture informative extrapolation.

\section*{Remaining work and source relationship}
Sampling uncertainty can be handled by enlarging (2) to a simultaneous confidence
region for all cell probabilities; any true compatible masses are then retained,
so optimizing (3) gives an outer confidence region for the whole identified set.
A useful procedure must also report width and positivity, since tiny conditional
cells can make coverage uninformative. The missing empirical ingredients remain
which policy responses are credible, how pilots represent the target population,
and how treatment versions and equilibrium spillovers are controlled.

\section{Completion Estimate}
60\%

This is a post hoc, heuristic estimate for the full catalogue target.
The attempt derives sharp universal-treatment bounds under a binary monotone instrument and an exact finite latent-type linear program for policy effects. Credible multidimensional policy responses, continuous-outcome error control, representative pilots, and equilibrium spillovers still require substantive identification.

\begin{thebibliography}{9}
\bibitem{source} M. Mogstad and A. Torgovitsky (2024),
\emph{Instrumental Variables with Unobserved Heterogeneity in Treatment Effects},
10 December author chapter. Sections 2.6, 4.3 and 4.7 provide context for
policy-specific IV and bounding methods.
\url{https://www.rfberlin.com/wp-content/uploads/2024/12/24030.pdf}.
The principal-stratum derivation and finite program above are standard-method
specializations with proofs, not claims of a new general IV solution.
\end{thebibliography}

\end{document}
